\documentclass{osmscience}

% Add bibliography file here (class does not add it)
\addbibresource{references.bib}

\title{Title of Abstract}
\author[1,*]{Firstname Lastname}
\author[1]{Firstname Lastname}
\author[2]{Firstname Lastname}
\affil[1]{
    Department/Lab/Unit, Name of Institution, City, Country\\
    \href{mailto:author1@institution1.com}{\underline{author1@institution1.com}}, 
    \quad
    \href{mailto:author2@institution1.com}{\underline{author2@institution1.com}}
}
\affil[2]{
    GIS Center, Florida International University, Miami, FL, USA\\
    \href{mailto:author@institution.edu}{\underline{author@institution.edu}}
}

% DOI Information 
% When uploading your abstract PDF to Zenodo you will be able 
% to generate a DOI - when this DOI is generated please copy 
% the DOI to this part of the template. 
\proceedingsDOI{10.1234/osmscience.2026.0001}

% to set the first page number for this paper in the proceedings
\proceedingsFirstPage{123}
\date{}

\begin{document}
\maketitle

\vspace{1em}

This is a paragraph. The content of your abstract should be inserted here as running text. 
Use the default "Normal text" style in this document. 
Notice that the first line of paragraphs are indented. 
Since this is an abstract, please \textbf{do not add headings for different sections} to ensure that all abstracts in the proceedings are formatted the same way. 
\textbf{Please avoid adding tables}.

This is a paragraph. 
The content of your abstract should be inserted here as running text. 
Use the default “Normal text” style in this document. 
Notice that the first line of paragraphs are indented. 
Since this is an abstract, please do not add headings for different sections to ensure that all abstracts in the proceedings are formatted the same way.

\textbf{\underline{References must be numbered}} and inserted in the text in the order they appear, e.g.~\cite{haklay2010good}. 
References can be cited together as \cite{barron2014comprehensive, de2007geospatial}. 
The \textbf{References list appearing at the end of the abstract} should be formatted according to the \textbf{APA style}. \textbf{You MUST not change the formatting of the references in this template as each entry in the OSMScience proceedings is expected to use the same style.}
For details/guidance, see: 
(i) the examples below, 
(ii) \url{http://www.bibme.org/citation-guide/apa}, 
(iii) the SotM proceedings of the previous years (2019: \url{https://zenodo.org/record/3405431#.YLzyK-dS9PY}, 2020: \url{https://zenodo.org/record/3928675#.YLzyBudS9PY}, 2021: \url{https://zenodo.org/record/5116434#.YrStenXMI88}). 
Notice that the first line of paragraphs are indented. 
The content of your abstract should be inserted here as running text. 
Use the default “Normal text” style in this document. 
Notice that the first line of paragraphs are indented. 
Since this is an abstract, please do not add headings for different sections to ensure that all abstracts in the proceedings are formatted the same way. 
This is a paragraph. 
The content of your abstract should be inserted here as running text. 
Use the default “Normal text” style in this document.

\textbf{The abstract cannot exceed 4 pages (including 1 figure max. and the References list) using this template.}
This rule will be strictly enforced to ensure the integrity of the proceedings. 
This is a paragraph. 
The content of your abstract should be inserted here as running text. 
Use the default “Normal text” style in this document. 
Notice that the first line of paragraphs are indented. 
Since this is an abstract, please do not add headings for different sections to ensure that all abstracts in the proceedings are formatted the same way. 
This is a paragraph. 
The content of your abstract should be inserted here as running text. 
Use the default “Normal text” style in this document. 
Notice that the first line of paragraphs are indented. 
Since this is an abstract, please do not add headings for different sections to ensure that all abstracts in the proceedings are formatted the same way.

\textbf{The abstract can include a maximum of 1 figure}, which shall always be referred to as \figref{fig:logo}. 
The figure shall be horizontally centered and shall have a caption, also centered and ending will a full-stop (see example below). 
\textbf{Please make sure that the figure is introduced in the text before appearing}, for example after a statement about its content (see \figref{fig:logo}).

\begin{figure}[ht]
    \centering
    \includegraphics[width=0.4\textwidth]{Openstreetmap_logo}
    \caption{This is the text of the caption and shows the logo of our beloved OSM.}
    \label{fig:logo}
\end{figure}

This is a paragraph. 
The content of your abstract should be inserted here as running text. 
Use the default “Normal text” style in this document. 
Notice that the first line of paragraphs are indented. 
Since this is an abstract, please do not add headings for different sections to ensure that all abstracts in the proceedings are formatted the same way. 
This is a paragraph. 
The content of your abstract should be inserted here as running text. 
Use the default “Normal text” style in this document. 
Notice that the first line of paragraphs are indented. 
Since this is an abstract, please do not add headings for different sections to ensure that all abstracts in the proceedings are formatted the same way.

For aesthetic reasons try to balance out the look of the last page as best as you can so that it wouldn’t contain only a few lines. 
In those cases, please shorten the text.
\printbibliography
 
\end{document}
